% ============================================================
\chapter{Th\'eor\`eme Central Limite}
\label{ch:tcl}
% ============================================================

Pourquoi la loi normale apparaît-elle partout --- dans la taille des individus, les erreurs de mesure, les rendements financiers, le bruit thermique~? La r\'eponse tient en un th\'eor\`eme~: le \emph{th\'eor\`eme central limite} (TCL). Abraham de Moivre en donne la premi\`ere version en 1733 pour les lancers de pi\`eces~; Laplace le g\'en\'eralise vers 1810~; Lindeberg (1922) et L\'evy (1925) \'etablissent la version d\'efinitive sous des hypoth\`eses minimales. Le TCL affirme que la somme normalis\'ee de variables al\'eatoires ind\'ependantes et identiquement distribu\'ees, quelle que soit leur loi commune (pourvu que la variance soit finie), converge en distribution vers la loi normale. C'est le r\'esultat le plus c\'el\`ebre --- et le plus utilis\'e --- de toute la th\'eorie des probabilit\'es.

\begin{intuition}
Le th\'eor\`eme central limite (TCL) est sans doute le r\'esultat le plus c\'el\`ebre
de la th\'eorie des probabilit\'es. Il affirme que la somme normalis\'ee de variables
al\'eatoires i.i.d.\ converge en loi vers la loi normale, quelle que soit la loi
commune des variables, pourvu que la variance soit finie.
\end{intuition}

% ------------------------------------------------------------
\section{\'Enonc\'e du th\'eor\`eme}
% ------------------------------------------------------------

\begin{theoreme}[Th\'eor\`eme central limite]
\label{thm:tcl}
Soit $(X_n)_{n \geq 1}$ une suite i.i.d.\ avec $\E[X_1] = \mu$ et
$\Var(X_1) = \sigma^2 \in (0, \infty)$. Posons $S_n = \sum_{i=1}^n X_i$.
Alors~:
\[
    Z_n = \frac{S_n - n\mu}{\sigma\sqrt{n}}
    = \frac{\bar{X}_n - \mu}{\sigma/\sqrt{n}}
    \xrightarrow{\mathcal{L}} \mathcal{N}(0, 1).
\]
Autrement dit, pour tout $x \in \R$~:
\[
    \lim_{n \to \infty} \PP(Z_n \leq x) = \Phi(x)
    = \frac{1}{\sqrt{2\pi}} \int_{-\infty}^{x} e^{-t^2/2}\, dt.
\]
\end{theoreme}

\begin{intuition}[title=\textbf{Pourquoi la loi normale est-elle omnipresente~?}]
Le TCL explique pourquoi tant de ph\'enom\`enes naturels suivent approximativement
une loi normale~: toute grandeur r\'esultant de la somme de nombreuses petites
contributions ind\'ependantes sera approximativement gaussienne.
\end{intuition}

% ------------------------------------------------------------
\section{D\'emonstration par les fonctions caract\'eristiques}
% ------------------------------------------------------------

Nous renvoyons au chapitre~11 pour la th\'eorie compl\`ete des fonctions
caract\'eristiques. Ici, nous esquissons la preuve en admettant le th\'eor\`eme
de L\'evy.

\begin{proof}[Preuve du TCL]
Sans perte de g\'en\'eralit\'e, posons $\mu = 0$ (sinon remplacer $X_i$ par
$X_i - \mu$). Soit $\varphi(t) = \E[e^{itX_1}]$ la fonction caract\'eristique
commune des $X_i$.

Par hypoth\`ese, $\varphi(0) = 1$, $\varphi'(0) = i\E[X_1] = 0$, et
$\varphi''(0) = i^2\E[X_1^2] = -\sigma^2$. Le d\'eveloppement de Taylor donne~:
\[
    \varphi(t) = 1 - \frac{\sigma^2 t^2}{2} + o(t^2) \quad (t \to 0).
\]

La fonction caract\'eristique de $Z_n = S_n/(\sigma\sqrt{n})$ est~:
\begin{align*}
    \varphi_{Z_n}(t) &= \E\!\left[\exp\!\left(it \cdot \frac{S_n}{\sigma\sqrt{n}}\right)\right]
    = \left[\varphi\!\left(\frac{t}{\sigma\sqrt{n}}\right)\right]^n \\
    &= \left[1 - \frac{\sigma^2}{2}\cdot\frac{t^2}{\sigma^2 n}
       + o\!\left(\frac{1}{n}\right)\right]^n
    = \left[1 - \frac{t^2}{2n} + o\!\left(\frac{1}{n}\right)\right]^n.
\end{align*}
En utilisant $\lim_{n \to \infty}(1 + a/n)^n = e^a$~:
\[
    \varphi_{Z_n}(t) \xrightarrow[n \to \infty]{} e^{-t^2/2}.
\]
Or $e^{-t^2/2}$ est la fonction caract\'eristique de $\mathcal{N}(0,1)$.
Par le th\'eor\`eme de continuit\'e de L\'evy, $Z_n \xrightarrow{\mathcal{L}}
\mathcal{N}(0,1)$.
\end{proof}

% ------------------------------------------------------------
\section{Applications pratiques}
% ------------------------------------------------------------

\subsection{Approximation de la loi binomiale}

\begin{corollaire}[Th\'eor\`eme de Moivre--Laplace]
Si $S_n \sim \text{Bin}(n, p)$ avec $0 < p < 1$, alors~:
\[
    \frac{S_n - np}{\sqrt{np(1-p)}} \xrightarrow{\mathcal{L}} \mathcal{N}(0,1).
\]
\end{corollaire}

\begin{exemple}
On lance une pi\`ece \'equilibr\'ee 100~fois. Quelle est la probabilit\'e
d'obtenir entre 45 et 55~Pile~?

Avec $S_{100} \sim \text{Bin}(100, 0{,}5)$, $\mu = 50$, $\sigma = 5$~:
\begin{align*}
    \PP(45 \leq S_{100} \leq 55) &= \PP\!\left(\frac{45 - 50}{5} \leq Z \leq \frac{55-50}{5}\right)
    \approx \PP(-1 \leq Z \leq 1) \\
    &= 2\Phi(1) - 1 \approx 0{,}6827.
\end{align*}
\end{exemple}

\begin{attention}[title=\textbf{Correction de continuit\'e}]
Pour am\'eliorer l'approximation, on utilise la \textbf{correction de continuit\'e}~:
\[
    \PP(a \leq S_n \leq b) \approx \Phi\!\left(\frac{b + 0{,}5 - np}{\sqrt{np(1-p)}}\right)
    - \Phi\!\left(\frac{a - 0{,}5 - np}{\sqrt{np(1-p)}}\right).
\]
\end{attention}

\subsection{Intervalles de confiance}

\begin{corollaire}[Intervalle de confiance pour la moyenne]
Si $(X_i)$ sont i.i.d.\ avec $\E[X_1] = \mu$ et $\Var(X_1) = \sigma^2$ connus,
un \textbf{intervalle de confiance \`a 95\%} pour $\mu$ est~:
\[
    \left[\bar{X}_n - 1{,}96\frac{\sigma}{\sqrt{n}},\;
    \bar{X}_n + 1{,}96\frac{\sigma}{\sqrt{n}}\right].
\]
\end{corollaire}

\begin{formules}[title=\textbf{Quantiles courants pour les IC}]
\begin{center}
\renewcommand{\arraystretch}{1.3}
\begin{tabular}{ccc}
    \toprule
    Niveau de confiance & $\alpha$ & $z_{\alpha/2}$ \\
    \midrule
    90\% & 0{,}10 & 1{,}645 \\
    95\% & 0{,}05 & 1{,}960 \\
    99\% & 0{,}01 & 2{,}576 \\
    \bottomrule
\end{tabular}
\end{center}
\end{formules}

\subsection{Approximation de la loi de Poisson}

\begin{corollaire}
Si $X \sim \text{Poisson}(\lambda)$ avec $\lambda$ grand, alors
$(X - \lambda)/\sqrt{\lambda} \xrightarrow{\mathcal{L}} \mathcal{N}(0,1)$.
\end{corollaire}

% ------------------------------------------------------------
\section{TCL multidimensionnel}
% ------------------------------------------------------------

\begin{theoreme}[TCL multidimensionnel]
Soit $(X_n)$ une suite i.i.d.\ de vecteurs al\'eatoires dans $\R^d$, avec
$\E[X_1] = \boldsymbol{\mu}$ et $\Cov(X_1) = \Sigma$. Alors~:
\[
    \sqrt{n}\,(\bar{X}_n - \boldsymbol{\mu})
    \xrightarrow{\mathcal{L}} \mathcal{N}_d(0, \Sigma).
\]
\end{theoreme}

% ------------------------------------------------------------
\section{M\'ethode delta}
% ------------------------------------------------------------

\begin{theoreme}[M\'ethode delta]
\label{thm:delta}
Si $\sqrt{n}(\bar{X}_n - \mu) \xrightarrow{\mathcal{L}} \mathcal{N}(0, \sigma^2)$
et si $g$ est d\'erivable en $\mu$ avec $g'(\mu) \neq 0$, alors~:
\[
    \sqrt{n}(g(\bar{X}_n) - g(\mu))
    \xrightarrow{\mathcal{L}} \mathcal{N}(0, \sigma^2 [g'(\mu)]^2).
\]
\end{theoreme}

\begin{exemple}
Si $\bar{X}_n \to p$ et $g(x) = x(1-x)$, alors
$g'(p) = 1-2p$ et~:
\[
    \sqrt{n}(\bar{X}_n(1-\bar{X}_n) - p(1-p))
    \xrightarrow{\mathcal{L}} \mathcal{N}(0, p(1-p)(1-2p)^2).
\]
\end{exemple}

% ------------------------------------------------------------
\section{Vitesse de convergence : th\'eor\`eme de Berry--Esseen}
% ------------------------------------------------------------

\begin{theoreme}[Berry--Esseen]
Sous les hypoth\`eses du TCL, si $\E[\abs{X_1}^3] < \infty$~:
\[
    \sup_{x \in \R} \abs{\PP(Z_n \leq x) - \Phi(x)}
    \leq \frac{C\, \E[\abs{X_1 - \mu}^3]}{\sigma^3 \sqrt{n}},
\]
o\`u $C$ est une constante universelle ($C \leq 0{,}4748$).
\end{theoreme}

\begin{remarque}
Berry--Esseen montre que la convergence dans le TCL est de l'ordre de $1/\sqrt{n}$.
Cela justifie l'utilisation de l'approximation normale d\`es que $n$ est
mod\'er\'ement grand (typiquement $n \geq 30$).
\end{remarque}

% ------------------------------------------------------------
\section{Exercices}
% ------------------------------------------------------------

\begin{exercice}
Soit $(X_i)$ i.i.d.\ avec $\E[X_1] = 2$ et $\Var(X_1) = 9$. Donner une
approximation de $\PP(S_{100} > 220)$.
\end{exercice}

\begin{exercice}
Un d\'e \'equilibr\'e est lanc\'e 360~fois. Approximer la probabilit\'e que la
somme d\'epasse 1300.
\end{exercice}

\begin{exercice}
Soit $S_n \sim \text{Bin}(400, 0{,}4)$. Calculer approximativement
$\PP(S_n \geq 180)$ avec et sans correction de continuit\'e.
\end{exercice}

\begin{exercice}
On mesure une grandeur physique $n = 50$ fois. Les mesures sont i.i.d.\ de
variance $\sigma^2 = 4$. Quel est l'intervalle de confiance \`a 99\% pour la
moyenne~?
\end{exercice}

\begin{exercice}
Combien de sondages faut-il r\'ealiser pour estimer une proportion $p$ \`a 2\%
pr\`es avec une confiance de 95\%~? (On ne conna\^it pas $p$.)
\end{exercice}

\begin{exercice}
En utilisant la m\'ethode delta, donner la loi limite de $\sqrt{n}(\ln\bar{X}_n - \ln\mu)$
lorsque les $X_i$ sont i.i.d.\ positives avec $\E[X_1] = \mu > 0$ et
$\Var(X_1) = \sigma^2$.
\end{exercice}

\begin{exercice}
V\'erifier le TCL pour $X_i \sim \text{Exp}(\lambda)$ en calculant directement
la fonction caract\'eristique de $Z_n$ et en montrant qu'elle converge vers
$e^{-t^2/2}$.
\end{exercice}

\begin{exercice}
Soit $(X_i)$ i.i.d.\ $\text{Poisson}(4)$. Approximer $\PP(\sum_{i=1}^{100} X_i \leq 380)$.
\end{exercice}
